% Overture Combinatorial Testing Guidelines
\documentclass{overturerepchap}
\usepackage{url}
\usepackage{graphics}
\usepackage{times}
\usepackage{listings}
\usepackage{color}
\usepackage{graphicx}
\usepackage{latexsym}
\usepackage{longtable} % ,multirow
\usepackage{vdmsl-2e}
\usepackage{makeidx}
\usepackage{hyperref}
\usepackage{ifthen}
\usepackage{makeidx}
\usepackage{fancyhdr}

\newcommand{\nonstandard}[1]{}
\newcommand{\vdmtools}{VDMTools}
\newcommand{\vdmstyle}[1]{\texttt{#1}}
\newcounter{exerciseno}
\newcommand{\vdmsl}{VDM-SL}
\newcommand{\vdmpp}{VDM++}
\newcommand{\Index}[1]{#1\index{#1}}

%\newcommand{\Exercise}[1]{%
%    \textbf{Exercise \thechapter.\theexerciseno}
%   \refstepcounter{exerciseno} #1 $\Box$\\ }%}
%\newcommand{\initexercise}{\setcounter{exerciseno}{0}}
%\newcounter{exerciseno}
\newcommand{\MYEQUIV}{$\equiv$}
\newenvironment{TypeSemantics}{\begin{longtable}[r]{|p{3.5cm}|p{9cm}|}\hline%
  Operator Name & Semantics Description \\ \hline\hline \endhead}%
  {\hline\end{longtable}}

\newcommand\thebookexercise{\thechapter.\arabic{exerciseno}}
\newenvironment{myexercise}{\par
  \refstepcounter{exerciseno}%
  \indent\textbf{Exercise\ \thebookexercise}\enskip}{$\Box$\\
}
\newenvironment{myhardexercise}{\par
  \refstepcounter{exerciseno}%
  \indent\textbf{Exercise\ \thebookexercise $\star$}\enskip}{$\Box$\\
}
\newcommand{\initexercise}{\setcounter{exerciseno}{0}}
%\newenvironment{mysolution}{
%} %% This will be replaced by a perl script extracting the solutions
  %% and inserting them automatically into the solutions chapter.
\newcommand{\insertfig}[4]{ % Filename, epsheight, epswidth, caption,  label
\begin{figure}[htb]
\begin{center}
\includegraphics[width=#2]{#1}
\end{center}
\caption{{\em #3}} #4
\end{figure}
}

\newcommand{\insertfignw}[3]{ % Filename, caption,  label
\begin{figure}[h]
\begin{center}
\includegraphics{#1}
\end{center}
\caption{{\em #2}} #3
\end{figure}
}

\newcommand{\insertfignwrot}[3]{ % Filename, caption,  label
\begin{figure}[h]
\begin{center}
\scalebox{1}%
{\rotatebox{270}{%
\includegraphics{#1}}}
\end{center}
\caption{{\em #2}} #3
\end{figure}
}

%\newcommand{\insertcommentedvdm*}[2]{}
  %% This macro is identified by perl script which will move the
  %% parameter to the solutions chapter.
% definition of VDM++, JavaCC, JJTree, JTB, ANTLR and SableCC for listings
\lstdefinelanguage{VDM++}
  {morekeywords={act, active, fin, req, waiting, abs, all, allsuper, always, and, answer, 
     assumption, async, atomic, be, bool, by, card, cases, char, class, comp, compose, conc, cycles,
     dcl, def, definitions, del, dinter, div, dlmodule, do, dom, dunion, duration, effect, elems, else, elseif, end,
     eq, error, errs, exists, exists1, exit, exports, ext, floor, for, forall, from, functions, 
     general, hd, if, imports, in, inds, infer, init, inmap, input, instance, int, inter, inv, inverse, iota, is, 
     isofbaseclass, isofclass, inv, inverse, lambda, len, let, map, measure, mu,
     mutex, mod, module, nat, nat1, new, merge, 
     munion, not, of, operations, or, ord, others, per, periodic, post, power, pre, pref, 
     private, protected, public, qsync, rd, responsibility, return, reverse,  
     sameclass, parameters, psubset, pure, rem, renamed, rng, sel, self, seq, seq1, set, skip, specified, st, 
     start, startlist, state, static, stop, stoplist, sporadic, subclass, subset, subtrace, sync, system, then, thread, 
     threadid, time, tixe, tl, to, token, traces, trap, types, undefined,
     union, uselib, using, values, 
     variables, while, with, wr, yet, RESULT, false, true, nil, periodic pref, rat, real},
   %keywordsprefix=mk\_,
   %keywordsprefix=a\_,
   %keywordsprefix=t\_,
   %keywordsprefix=w\_,
   sensitive,
   morecomment=[l]--,
   morestring=[b]",
   morestring=[b]',
  }[keywords,comments,strings]
\lstdefinelanguage{JavaCC}
  {morekeywords={options, PARSER\_BEGIN, PARSER\_END, SKIP, TOKEN},
   sensitive=false,
  }[keywords]

% define the layout for listings
\lstdefinestyle{tool}{basicstyle=\ttfamily,
                         frame=trBL,
			 showstringspaces=false,
			 frameround=ffff,
			 framexleftmargin=0mm,
			 framexrightmargin=0mm}
\lstdefinestyle{mystyle}{basicstyle=\ttfamily,
                         frame=trBL,
%                         numbers=left,
%			 gobble=0,
			 showstringspaces=false,
%			 linewidth=\textwidth,
			 frameround=fttt,
			 aboveskip=5mm,
			 belowskip=5mm,
			 framexleftmargin=0mm,
			 framexrightmargin=0mm}
%\lstdefinestyle{mystyle}{basicstyle=\sffamily\small,
%			 frame=tb,
%                         numbers=left,
%			 gobble=0,
%			 showstringspaces=false,
%			 linewidth=345pt,
%			 frameround=ffff,
%			 framexleftmargin=8mm,
%			 framexrightmargin=8mm,
%			 framextopmargin=1mm,
%			 framexbottommargin=1mm,
%			 aboveskip=7mm,
%			 belowskip=5mm,
%			 xleftmargin=10mm,}

\lstset{style=mystyle}
\lstset{language=VDM++}
\lstset{alsolanguage=Java}
% The command below enables you to escape into normal LaTeX mode inside your 
% VDM chunks by starting with a `!\UTF{FFFD}\UTF{FFFD}? character and ending with a `!\UTF{FFFD}\UTF{FFFD}?
\lstset{escapeinside=!!}

%\newcommand{\kw}[1]{{\tt #1}}

%%%%%%%%%%%%%%%%%% Commands for bibtex %%%%%%%%%%%%%%%%%%%%%%%
%************************************************************************
%                                                                       *
%       Bibliography and Terminology supporting commands                *
%                                                                       *
%************************************************************************

\newcommand{\bthisbibliography}[1]{\chapter*{References}%
   \begin {list} {}%
     {\settowidth {\labelwidth} {[#1]XX}%
      \setlength {\leftmargin} {\labelwidth}%
      \addtolength{\leftmargin} {\labelsep}%
      \setlength {\parsep} {1ex}%
      \setlength {\itemsep} {2ex}%
     }
  }
\newcommand{\ethisbibliography}{\end{list}}
\newcommand{\refitem}[2]
  {\bibitem[#1]{#2}}

% Requirements environment
\newenvironment{reqs}{%
\begin{enumerate}
%\renewcommand{\labelenumi}{\textbf{R\theenumi}}
\renewcommand{\theenumi}{\textbf{R\arabic{enumi}}}
}{%
\end{enumerate}}

%\newcommand{\Exercise}[1]{%
%    \textbf{Exercise \thechapter.\theexerciseno}
%   \refstepcounter{exerciseno} #1 $\Box$\\ }%}
%\newcommand{\initexercise}{\setcounter{exerciseno}{0}}
\newcommand{\Lit}[1]{`{\tt #1}\Quote}
\newcommand{\Rule}[2]{
  \begin{quote}\begin{tabbing}
    #1\index{#1}\ \ \= = \ \ \= #2  ; %    Adds production rule to index

  \end{tabbing}\end{quote}
  }
\newcommand{\RuleTarget}[1]{\hypertarget{rule:#1}{}}
\newcommand{\Ruledef}[2]
{
  \RuleTarget{#1}\Rule{#1}{#2}%
  }
\newcommand{\Ruleref}[1]{
  \hyperlink{rule:#1}{#1}}
\newcommand{\SeqPt}[1]{\{\ #1\ \}}
\newcommand{\lfeed}{\\ \> \>}
\newcommand{\dsepl}{\ $|$\ }
\newcommand{\dsep}{\\ \> $|$ \>}
\newcommand{\Lop}[1]{`{\bf\ttfamily #1}\Quote}
\newcommand{\blankline}{\vspace{\baselineskip}}
\newcommand{\Brack}[1]{(\ #1\ )}
\newcommand{\nmk}{\footnotemark}
\newcommand{\ntext}[1]{\footnotetext{{\bf Note: } #1}}
\newlength{\kwlen}
\newcommand{\Keyw}[1]{\settowidth{\kwlen}{\tt
    #1}\makebox[\kwlen][l]{{\bf\ttfamnily #1}}}
\newcommand{\keyw}[1]{{\bf\ttfamily #1}}
\newcommand{\id}[1]{{\tt #1}}
\newcommand{\metaiv}[1]{\begin{lstlisting}\input{#1}\end{lstlisting}}

\newcommand{\OptPt}[1]{[\ #1\ ]}
\newcommand{\MAP}[2]{\keyw{map }#1\keyw{ to }#2}
\newcommand{\INMAP}[2]{\keyw{inmap }#1\keyw{ to }#2}
\newcommand{\SEQ}[1]{\keyw{seq of }#1}
\newcommand{\NSEQ}[1]{\keyw{seq1 of }#1}
\newcommand{\SET}[1]{\keyw{set of }#1}
\newcommand{\SETONE}[1]{\keyw{set1 of }#1}
\newcommand{\PROD}[2]{#1 * #2}
\newcommand{\TO}[2]{#1 $\To$ #2}
\newcommand{\FUN}[2]{#1 $\To$ #2}
\newcommand{\PUBLIC}{\ifthenelse{\boolean{VDMpp}}{public\mbox{}}{\mbox{}}}
\newcommand{\PRIVATE}{\ifthenelse{\boolean{VDMpp}}{private}{\mbox{}}}
\newcommand{\PROTECTED}{\ifthenelse{\boolean{VDMpp}}{protected}{\mbox{}}}

\newcommand{\experience}[1]{%
\begin{center}
\fbox{
\begin{minipage}[t]{.8\textwidth}
#1
\end{minipage}}
\end{center}}


\pagestyle{fancy}
\fancyhead{}
\fancyhead[LO]{\leftmark}
\fancyhead[RE]{The VDM Combinatorial Testing Guidelines}
\fancyhead[RO,LE]{\resizebox{0.05\textwidth}{!}{\includegraphics{overture}}}
\fancyfoot[C]{\thepage}
\makeindex

\begin{document}

\title{Guidelines for using VDM Combinatorial Testing Features}
\author{Nick Battle\\
            Peter Gorm Larsen}

\date{December 2020}
\reportno{TR-005}

\pagenumbering{roman}
\maketitle


{\textbf{Document history}}

\begin{tabular}{|l|l|l|l|l|}\hline
Month   & Year & Version & Version of Overture.exe & Comment \\ \hline
December   & 2010 &    0.1     & 3.0.2   & Initial version\\ \hline
\end{tabular}

%\addtocounter{page}{2}
\tableofcontents
\newpage
\mbox{}
\newpage
\pagenumbering{arabic}
\setcounter{page}{1}

\chapter{Introduction}

This manual is a complete guide the combinatorial testing of VDM models. It
assumes the reader has no prior knowledge of combinatorial testing, but a
working knowledge of VDM, in particular, the VDM++ and VDM-SL dialects.

\section{What is Combinatorial Testing?}

Creating a comprehensive set of tests for VDM specifications can be a time
consuming process. To try to make the generation of test cases simpler, Overture
provides a VDM language extension (for all dialects) called Combinatorial
Testing \cite{Nie&11,Larsen&10c}.

In general, specifications are tested to verify that certain properties or
behaviours are met, as specified by constraints in the specification and
validation conjectures in the tests.

The simplest way to test a specification is to write ad-hoc tests, starting from
a known system state and proceeding with a sequence of operation calls that
should move to a new state or produce some particular result or error response.
The problem with this kind of testing is that it can be very laborious to produce
the number of tests needed to cover the complete system behaviour. It is also
expensive to maintain a large test suite as the specification evolves.

The most complete way to test a specification is to produce a formal
mathematical proof that it will never violate its constraints, and always meet
its validation conjectures if presented with a legal sequence of operation
calls. This provides the highest level of confidence in the correctness of a
specification, but it can be unrealistic to produce a complete formal proof for
complex specifications, even with tool support \cite{Paulson97,Bicarregui&94}.

Model checking provides an approach to formal testing that is considerably
better than ad-hoc testing but not as complete as formal proof \cite{Clarke&99}. This approach
uses a formal specification of the system properties desired, often written in a
temporal calculus, and the model checker symbolically executes the
specification searching for execution paths that violate the constraints. Since
the execution is symbolic, extremely large state spaces can be searched
(billions of cases is not uncommon), and failed cases can produce a ``counter
example'' that demonstrates the failure. This is a very powerful technique, but
in practice, realistic specifications often produce a state space explosion
that is too great for model checkers.

Combinatorial testing is an approach that is far more powerful than ad-hoc
testing, but not as complete as model checking. Tests are produced automatically
from ``\texttt{\textbf{traces}}'' that are relatively simple to define. The approach allows
specifications to be tested with perhaps millions of test cases, but cannot
guarantee to catch every corner case in the way that a model checker can.
Therefore the technique is useful for specifications that are too complex for
model checking or formal proof.

A combinatorial trace is a pattern that describes the construction of
argument values and the sequences of operation calls that will exercise the
specification. A specification may contain several traces, each designed to test
a particular aspect. Traces are automatically expanded into a (potentially
large) number of tests, each of which is a particular sequence of operation
calls and argument values. The execution of tests is performed
automatically, starting each in a known state; a test is considered to pass if
it does not violate the specification's constraints, or the test's validation
conjectures. Individual failed tests can be executed in isolation to find out
why they failed, which is similar to a model checker's counter example. 

\section{The Structure of this Document}

The document is intended to be read sequentially, but readers who are familiar
with the basics of VDM traces can skip Chapter \ref{chap:traces} and go straight
to the patterns in Chapter \ref{chap:patterns}.

\begin{itemize}
\item Chapter \ref{chap:traces}, \nameref{chap:traces} introduces the concept of
a combinatorial test and explains how to work with traces in VDM.
\item Chapter \ref{chap:patterns}, \nameref{chap:patterns} looks at common ways
of using traces that are useful in many different situations.
\item Chapter \ref{chap:examples}, \nameref{chap:examples} uses two more
significant specifications to demonstrate typical usage of traces in real life.
\item Appendix \ref{chap:CTsyntax}, \nameref{chap:CTsyntax} defines the formal
grammar of traces.
\item Appendix \ref{chap:OvertureShots}, \nameref{chap:OvertureShots} includes
screen shots of the Combinatorial Testing perspective in Overture.
\item Appendix \ref{chap:listings}, \nameref{chap:listings} contains the full
listings of the models explored in Chapter \ref{chap:examples}.
\end{itemize}

\chapter{Working with Traces}
\label{chap:traces}

\section{Basic Trace Constructs}
Combinatorial tests are embedded within a VDM specification using a section
called ``\texttt{\textbf{traces}}''. Typically, one or more traces are added to a separate class
or module that is intended for testing rather than the main specification,
though you can add traces to any class you wish. In this chapter, we will use
the example classes below:

\small
\begin{lstlisting}
class Counter
instance variables
    total:int := 0;

operations
    public inc: () ==> int
    inc() == ( total := total + 1; return total; )

    public dec: () ==> int
    dec() == ( total := total - 1; return total; )

end Counter

class Tester
instance variables
    obj:Counter := new Counter();

traces
    T1: obj.inc();

end Tester
\end{lstlisting}
\normalsize

Notice that there are two classes, \texttt{Counter} and \texttt{Tester}. The \texttt{Counter} class defines a
simple operation that increments and decrements a total state value that is
initially zero. The \texttt{Tester} class creates an instance of Counter and defines a
single trace called \texttt{T1}. Trace names are simple identifiers, optionally separated
by slashes (e.g.\ \texttt{item456/interface/all}). This example is the simplest
trace possible and indicates that the trace should expand to a single test that
just calls \texttt{obj.inc()}.

This trace can either be executed in Overture in the Combinatorial Testing
perspective (see Appendix \ref{chap:OvertureShots}), or it can be executed from
the command line using the \texttt{runtrace} command. The command line output is
illustrated here for simplicity:

\lstset{style=tool,language=}
\begin{lstlisting}[escapechar=@]
> runtrace Tester`T1
Generated 1 tests in 0.004 secs. 
Test 1 = obj.inc()
Result = [1, PASSED]
Executed in 0.007 secs. 
All tests passed
\end{lstlisting}
\lstset{style=mystyle}
\lstset{language=VDM++}

The first line of output indicates that one test has been generated from the
trace. With more complex examples, this generation could expand to thousands or
millions of tests, and consequently it may take a few seconds.

The next line of the output describes the test that was generated. Note that
this is called ``Test 1'', and consists of a single call to \texttt{obj.inc()}.

The line below the test gives the result of executing that test. There is a
single return value from the call to \texttt{obj.inc()}, 1, which is listed
along with the word ``PASSED'' that indicates that there were no constraint violations
in the test execution.

Lastly the time taken to execute all of the tests is given, and an indication of
whether any tests failed.

The reason that this trace only expands to a single test is that the trace, when
considered as a pattern, only matches a single operation call. But if we change
the trace to the following:

\small
\begin{lstlisting}
traces
    T1: obj.inc() | obj.dec();
\end{lstlisting}
\normalsize

The trace is now saying that it would match either a call to \texttt{obj.inc()}
or a call to \texttt{obj.dec()}. Therefore the test expansion produces the
following:

\lstset{style=tool,language=}
\begin{lstlisting}[escapechar=@]
> runtrace Tester`T1
Generated 2 tests
Test 1 = obj.inc()
Result = [1, PASSED]
Test 2 = obj.dec()
Result = [-1, PASSED]
Executed in 0.033 secs. 
All tests passed
\end{lstlisting}
\lstset{style=mystyle}
\lstset{language=VDM++}

This time two tests are generated. The first calls \texttt{obj.inc()}, the
second \texttt{obj.dec()}. Notice that the decrement test is completely separate
from the increment test. It produces -1 as its result, because the \texttt{Counter}
object is re-created for each test. It does not decrement the counter back to
zero after the first test incremented it.

If we want to test an increment followed by a decrement, that would be expressed
using a semi-colon separator:

\small
\begin{lstlisting}
traces
    T1: obj.inc(); obj.dec();
\end{lstlisting}
\normalsize

This produces the output:

\small
\lstset{style=tool,language=}
\begin{lstlisting}[escapechar=@]
Generated 1 tests
Test 1 = obj.inc(); obj.dec()
Result = [1, 0, PASSED]
Executed in 0.027 secs. 
All tests passed
\end{lstlisting}
\lstset{style=mystyle}
\lstset{language=VDM++}
\normalsize

This generates a single test again, but you can see that the test involves two
calls and that they return 1 and 0, respectively. So this time the second call
is operating on the same object instance as the first.

If the increment and decrement operations are independent, it makes sense to
test calls to them in either order, which would be expressed as:

\small
\begin{lstlisting}
traces
    T1: || ( obj.inc(), obj.dec() );
\end{lstlisting}
\normalsize

Notice that the separator has changed to a comma. That produces the output:

\small
\lstset{style=tool,language=}
\begin{lstlisting}[escapechar=@]
Generated 2 tests
Test 1 = obj.inc(); obj.dec()
Result = [1, 0, PASSED]
Test 2 = obj.dec(); obj.inc()
Result = [-1, 0, PASSED]
Executed in 0.029 secs. 
All tests passed
\end{lstlisting}
\lstset{style=mystyle}
\lstset{language=VDM++}
\normalsize

So now both orderings of the two calls are produced. This is because there are
two orderings that match the pattern \texttt{|| ( \ldots, \ldots )}. This
particular trace construct naturally expands to an arbitrary number of calls and
produces a test for every permutation of the calls in brackets.

But what if some tests are a pair of calls and some are not? If we want to make
a call optional, the \texttt{?} operator can be added to any operation call
(i.e.\ not just within \texttt{||} operators) to indicate that this will match
tests where the call is made and where it is not. For example:

\small
\begin{lstlisting}
traces
    T1: || ( obj.inc(), obj.dec()? );
\end{lstlisting}
\normalsize

This means that the decrement call is optional and so although it is included in
the orderings of the pair, it should also be absent in some cases. This
example produces the following:

\small
\lstset{style=tool,language=}
\begin{lstlisting}[escapechar=@]
Generated 4 tests
Test 1 = obj.inc(); skip
Result = [1, (), PASSED]
Test 2 = obj.inc(); obj.dec()
Result = [1, 0, PASSED]
Test 3 = skip; obj.inc()
Result = [(), 1, PASSED]
Test 4 = obj.dec(); obj.inc()
Result = [-1, 0, PASSED]
Executed in 0.043 secs. 
All tests passed
\end{lstlisting}
\lstset{style=mystyle}
\lstset{language=VDM++}
\normalsize

You see that the decrement call is sometimes present and sometimes replaced by
\texttt{\textbf{skip}}, which indicates the absence of an optional call. Notice also that
the \texttt{||} operator and the \texttt{?} operator work together to combine their
effects in this example, though \texttt{?} can be used for any operation call.

Along the same lines as \texttt{?}, it is possible to add \texttt{*} and
\texttt{+} operators to any call, which indicate that it should be called zero
or more times, and one or more times. The maximum number of times is a tool preset
value that defaults to 5, though it can be changed. So for example:

\small
\begin{lstlisting}
traces
    T1: obj.inc()*;
    T2: obj.dec()+;
\end{lstlisting}

\lstset{style=tool,language=}
\begin{lstlisting}[escapechar=@]
> runtrace Tester`T1
Generated 6 tests
Test 1 = skip
Result = [(), PASSED]
Test 2 = obj.inc()
Result = [1, PASSED]
Test 3 = obj.inc(); obj.inc()
Result = [1, 2, PASSED]
Test 4 = obj.inc(); obj.inc(); obj.inc()
Result = [1, 2, 3, PASSED]
Test 5 = obj.inc(); obj.inc(); obj.inc(); obj.inc()
Result = [1, 2, 3, 4, PASSED]
Test 6 = obj.inc(); obj.inc(); obj.inc(); obj.inc(); obj.inc()
Result = [1, 2, 3, 4, 5, PASSED]
Executed in 0.046 secs. 
All tests passed

> runtrace Tester`T2
Generated 5 tests
Test 1 = obj.dec()
Result = [-1, PASSED]
Test 2 = obj.dec(); obj.dec()
Result = [-1, -2, PASSED]
Test 3 = obj.dec(); obj.dec(); obj.dec()
Result = [-1, -2, -3, PASSED]
Test 4 = obj.dec(); obj.dec(); obj.dec(); obj.dec()
Result = [-1, -2, -3, -4, PASSED]
Test 5 = obj.dec(); obj.dec(); obj.dec(); obj.dec(); obj.dec()
Result = [-1, -2, -3, -4, -5, PASSED]
Executed in 0.042 secs. 
All tests passed
\end{lstlisting}
\lstset{style=mystyle}
\lstset{language=VDM++}
\normalsize

The important difference between these two is that \texttt{T1} includes an extra
\texttt{\textbf{skip}} case, whereas \texttt{T2} does not.

Lastly, it is possible to indicate a specific number of repetitions of a call or
a range of repetitions. For example:

\small
\begin{lstlisting}
traces
    T1: obj.inc(){3};
    T2: obj.dec(){2, 4};
\end{lstlisting}

\lstset{style=tool,language=}
\begin{lstlisting}[escapechar=@]
> runtrace Tester`T1
Generated 1 tests
Test 1 = obj.inc(); obj.inc(); obj.inc()
Result = [1, 2, 3, PASSED]
Executed in 0.026 secs. 
All tests passed

> runtrace Tester`T2
Generated 3 tests
Test 1 = obj.dec(); obj.dec()
Result = [-1, -2, PASSED]
Test 2 = obj.dec(); obj.dec(); obj.dec()
Result = [-1, -2, -3, PASSED]
Test 3 = obj.dec(); obj.dec(); obj.dec(); obj.dec()
Result = [-1, -2, -3, -4, PASSED]
Executed in 0.01 secs. 
All tests passed
\end{lstlisting}
\lstset{style=mystyle}
\lstset{language=VDM++}
\normalsize

The \texttt{T1} trace now produces a single test with precisely three repetitions, while
the \texttt{T2} trace gives three tests with 2, 3 and 4 repetitions respectively.

If you combine a \texttt{||} operator with a repetition, the result is to repeat
all of the possibilities of the permutation with the given number of repetitions.
For example:

\small
\begin{lstlisting}
traces
    T1: || ( obj.inc(), obj.dec() ) {2};
\end{lstlisting}

\lstset{style=tool,language=}
\begin{lstlisting}[escapechar=@]
> runtrace Tester`T1
Generated 4 tests
Test 1 = obj.inc(); obj.dec(); obj.inc(); obj.dec()
Result = [1, 0, 1, 0, PASSED]
Test 2 = obj.dec(); obj.inc(); obj.inc(); obj.dec()
Result = [-1, 0, 1, 0, PASSED]
Test 3 = obj.inc(); obj.dec(); obj.dec(); obj.inc()
Result = [1, 0, -1, 0, PASSED]
Test 4 = obj.dec(); obj.inc(); obj.dec(); obj.inc()
Result = [-1, 0, -1, 0, PASSED]
Executed in 0.038 secs. 
All tests passed
\end{lstlisting}
\lstset{style=mystyle}
\lstset{language=VDM++}
\normalsize

Here, the \texttt{||} operator produces (\texttt{inc}, \texttt{dec}) and (\texttt{dec}, \texttt{inc}); then the
repetition doubles this, but it doubles every combination of the two rather than
simply repeating each one twice.

\section{Using Variables}

So far, the trace examples have called operations that do not include any
arguments. Arguments can be passed as literals, but traces also provide the
means to define variables that can change value as tests are generated from a
trace.

If we overload the example increment and decrement operations with versions that
take an integer parameter, by which to change the counter, we can write traces
like this:

\small
\begin{lstlisting}
...
    public inc: int ==> int
    inc(i) == ( total := total + i; return total; );

    public dec: int ==> int
    dec(i) == ( total := total - i; return total; )

traces
    T1:
        let a in set {1, ..., 10} be st a mod 2 = 0 in
            obj.inc(a);
\end{lstlisting}

\lstset{style=tool,language=}
\begin{lstlisting}[escapechar=@]
> runtrace Tester`T1
Generated 5 tests
Test 1 = obj.inc(2)
Result = [2, PASSED]
Test 2 = obj.inc(4)
Result = [4, PASSED]
Test 3 = obj.inc(6)
Result = [6, PASSED]
Test 4 = obj.inc(8)
Result = [8, PASSED]
Test 5 = obj.inc(10)
Result = [10, PASSED]
Executed in 0.04 secs. 
All tests passed
\end{lstlisting}
\lstset{style=mystyle}
\lstset{language=VDM++}
\normalsize

In a standard VDM specification, the \texttt{\textbf{let}\ldots \textbf{be st}} expression would
choose an arbitrary element from the set that meets the \texttt{st} clause. But
in a trace context, this looseness is used as a pattern that expands to a test
covering each possible set value that would match. Notice that the tests list
the actual value of the argument passed, rather than the symbolic name, ``a''.

A trace can include multiple \texttt{\textbf{let}} clauses, but if these are nested, then
the trace expands to the \emph{combination} of the variables. For example:

\small
\begin{lstlisting}
traces
    T1:
        let a in set {1, 2, 3} in
            let b in set {4, 5, 6} in
                ( obj.inc(a); obj.dec(b) );
\end{lstlisting}

\lstset{style=tool,language=}
\begin{lstlisting}[escapechar=@]
> runtrace Tester`T1
Generated 9 tests
Test 1 = obj.inc(1); obj.dec(4)
Result = [1, -3, PASSED]
Test 2 = obj.inc(1); obj.dec(5)
Result = [1, -4, PASSED]
Test 3 = obj.inc(1); obj.dec(6)
Result = [1, -5, PASSED]
Test 4 = obj.inc(2); obj.dec(4)
Result = [2, -2, PASSED]
Test 5 = obj.inc(2); obj.dec(5)
Result = [2, -3, PASSED]
Test 6 = obj.inc(2); obj.dec(6)
Result = [2, -4, PASSED]
Test 7 = obj.inc(3); obj.dec(4)
Result = [3, -1, PASSED]
Test 8 = obj.inc(3); obj.dec(5)
Result = [3, -2, PASSED]
Test 9 = obj.inc(3); obj.dec(6)
Result = [3, -3, PASSED]
Executed in 0.066 secs. 
All tests passed
\end{lstlisting}
\lstset{style=mystyle}
\lstset{language=VDM++}
\normalsize

This example produces a test for every combination of ``\texttt{a}'' and ``\texttt{b}'' values,
which is therefore nine tests. The round brackets are needed around the pair of
operation calls because a call binds tightly to the \texttt{\textbf{let}}. Without the
brackets, you get the following scope error, referring to the ``a'' in the
second call to \texttt{obj.dec(a)}:

\small
\begin{lstlisting}
traces
    T1:
        let a in set {6, 7, 10} in
            obj.inc(a); obj.dec(a)
\end{lstlisting}
\scriptsize
\lstset{style=tool,language=}
\begin{lstlisting}[escapechar=@]
Error 3182: Name 'Tester`a' is not in scope in 'Tester' (example.vpp) at line 28:29
Type checked 2 classes in 0.12 secs. Found 1 type error
\end{lstlisting}
\lstset{style=mystyle}
\lstset{language=VDM++}
\normalsize

Note also that the variables defined are in scope
throughout the clauses below, so the ``\texttt{a}'' variable could be used to define the
set of ``\texttt{b}'' values:

\small
\begin{lstlisting}
traces
    T1:
        let a in set {1, 2, 3} in
            let b in set {a, ..., a + 2} in
                ( obj.inc(a); obj.dec(b) );
\end{lstlisting}

\lstset{style=tool,language=}
\begin{lstlisting}[escapechar=@]
> runtrace Tester`T1
Generated 9 tests
Test 1 = obj.inc(1); obj.dec(1)
Result = [1, 0, PASSED]
Test 2 = obj.inc(1); obj.dec(2)
Result = [1, -1, PASSED]
Test 3 = obj.inc(1); obj.dec(3)
Result = [1, -2, PASSED]
Test 4 = obj.inc(2); obj.dec(2)
Result = [2, 0, PASSED]
Test 5 = obj.inc(2); obj.dec(3)
Result = [2, -1, PASSED]
Test 6 = obj.inc(2); obj.dec(4)
Result = [2, -2, PASSED]
Test 7 = obj.inc(3); obj.dec(3)
Result = [3, 0, PASSED]
Test 8 = obj.inc(3); obj.dec(4)
Result = [3, -1, PASSED]
Test 9 = obj.inc(3); obj.dec(5)
Result = [3, -2, PASSED]
Executed in 0.059 secs. 
All tests passed
\end{lstlisting}
\lstset{style=mystyle}
\lstset{language=VDM++}
\normalsize

If two variables should take values from the same set of values, it is possible
to use a \texttt{multiple bind} in a trace, but not a \texttt{bind list}. For
example:

\small
\begin{lstlisting}
traces
    T1:
        let a, b in set {1, 2, 3} in
            ( obj.inc(a); obj.dec(b) );
\end{lstlisting}

\lstset{style=tool,language=}
\begin{lstlisting}[escapechar=@]
> runtrace Tester`T1
Generated 9 tests
Test 1 = obj.inc(1); obj.dec(1)
Result = [1, 0, PASSED]
Test 2 = obj.inc(2); obj.dec(1)
Result = [2, 1, PASSED]
Test 3 = obj.inc(3); obj.dec(1)
Result = [3, 2, PASSED]
Test 4 = obj.inc(1); obj.dec(2)
Result = [1, -1, PASSED]
Test 5 = obj.inc(2); obj.dec(2)
Result = [2, 0, PASSED]
Test 6 = obj.inc(3); obj.dec(2)
Result = [3, 1, PASSED]
Test 7 = obj.inc(1); obj.dec(3)
Result = [1, -2, PASSED]
Test 8 = obj.inc(2); obj.dec(3)
Result = [2, -1, PASSED]
Test 9 = obj.inc(3); obj.dec(3)
Result = [3, 0, PASSED]
Executed in 0.061 secs. 
All tests passed
\end{lstlisting}
\lstset{style=mystyle}
\lstset{language=VDM++}
\normalsize


As well as defining a variable value from a set, variables can be used to
simplify calculations that would otherwise have to be made in the arguments to
operation calls. These simpler \texttt{\textbf{let}} definitions do not increase the
number of tests generated from the trace, they just introduce new names in the
scope that follows. Multiple variable definitions can be declared in 
one \texttt{\textbf{let}} expression. For example:

\small
\begin{lstlisting}
traces
T1:
    let a in set {1, 2, 3} in
        let b = a + 1, c = a - 1 in
            ( obj.inc(b); obj.dec(c) );
\end{lstlisting}

\lstset{style=tool,language=}
\begin{lstlisting}[escapechar=@]
> runtrace Tester`T1
Generated 3 tests
Test 1 = obj.inc(2); obj.dec(0)
Result = [2, 2, PASSED]
Test 2 = obj.inc(3); obj.dec(1)
Result = [3, 2, PASSED]
Test 3 = obj.inc(4); obj.dec(2)
Result = [4, 2, PASSED]
Executed in 0.054 secs. 
All tests passed
\end{lstlisting}
\lstset{style=mystyle}
\lstset{language=VDM++}
\normalsize

If repetitions are added to a clause within a \texttt{\textbf{let}} body, they bind
tightly to the operation call rather than the entire \texttt{\textbf{let}} clause.
If you want to repeat the entire \texttt{\textbf{let}}, you have to
bracket the whole clause and add a repetition to that. For example:

\small
\begin{lstlisting}
traces
    T1:
        let a in set {1, 2, 3} in
            obj.inc(a){1, 2}
    T2:
        ( let a in set {1, 2, 3} in
            obj.inc(a) ){1, 2}
\end{lstlisting}

\lstset{style=tool,language=}
\begin{lstlisting}[escapechar=@]
> runtrace Tester`T1
Generated 6 tests
Test 1 = obj.inc(1)
Result = [1, PASSED]
Test 2 = obj.inc(1); obj.inc(1)
Result = [1, 2, PASSED]
Test 3 = obj.inc(2)
Result = [2, PASSED]
Test 4 = obj.inc(2); obj.inc(2)
Result = [2, 4, PASSED]
Test 5 = obj.inc(3)
Result = [3, PASSED]
Test 6 = obj.inc(3); obj.inc(3)
Result = [3, 6, PASSED]
Executed in 0.044 secs. 
All tests passed

> runtrace Tester`T2
Generated 12 tests
Test 1 = obj.inc(1)
Result = [1, PASSED]
Test 2 = obj.inc(2)
Result = [2, PASSED]
Test 3 = obj.inc(3)
Result = [3, PASSED]
Test 4 = obj.inc(1); obj.inc(1)
Result = [1, 2, PASSED]
Test 5 = obj.inc(2); obj.inc(1)
Result = [2, 3, PASSED]
Test 6 = obj.inc(3); obj.inc(1)
Result = [3, 4, PASSED]
Test 7 = obj.inc(1); obj.inc(2)
Result = [1, 3, PASSED]
Test 8 = obj.inc(2); obj.inc(2)
Result = [2, 4, PASSED]
Test 9 = obj.inc(3); obj.inc(2)
Result = [3, 5, PASSED]
Test 10 = obj.inc(1); obj.inc(3)
Result = [1, 4, PASSED]
Test 11 = obj.inc(2); obj.inc(3)
Result = [2, 5, PASSED]
Test 12 = obj.inc(3); obj.inc(3)
Result = [3, 6, PASSED]
Executed in 0.03 secs. 
All tests passed
\end{lstlisting}
\lstset{style=mystyle}
\lstset{language=VDM++}
\normalsize

The difference may seem subtle, but the effect is significant. T1 behaves like a
simple ``\{1, 2\}'' repetition for each of the \texttt{\textbf{let}} values, whereas T2
produces either one or two cases from the \emph{entire set} created by the \texttt{\textbf{let}}
clause.

Although the example above uses a \texttt{let <set bind>} expression, it is also
possible to use a \texttt{\textbf{let} <seq bind>} or \texttt{\textbf{let} <type bind>}. In a
trace context, the ordering that a sequence bind carries does not affect the
generation of tests; if the example had \texttt{\textbf{let} a \textbf{in seq} [1,2,3]}, the
same tests would be generated, and since tests are independent their order is not
meaningful. So sequence binds are not particularly useful in traces. However
type binds (of finite types) are a shorthand for ``all values of this type'',
which can be useful in some circumstances. This is covered later in Chapter
\ref{chap:patterns}.

\section{Tests with Errors}

The examples so far have only includes tests that \textbf{PASSED}. This means that they
completed the sequence of operation calls without violating any pre-conditions,
post-conditions, state invariants, type invariants, recursive measures or dynamic
type checks.

If a sequence of operations causes a post-condition failure, then it
is certain that there is a problem with the specification –- it should not be
possible to provoke a post-condition failure with a set of legal calls (ie. ones
which pass the pre-conditions and type invariants). On the other hand, if a 
sequence of operations violates a pre-condition, or a type or class invariant, 
then it is \emph{possible} that the specification has a problem, but it is also 
possible that the test itself is at fault (passing illegal values).

The combinatorial testing environment indicates the exit status of the test  in
the verdict returned in the last item of the results (all \texttt{PASSED} above). So if 
pre/post/invariant conditions are violated during a test, this may be set to 
\texttt{FAILED} or \texttt{INDETERMINATE}\footnote{The tools sometimes call this \texttt{INCONCLUSIVE}}. If
a test fails, then any subsequent test which starts \emph{with the same
sequence of calls} as the failed sequence will also fail. These tests are
filtered out of the remaining test sequence automatically, and not executed.

For example, if we introduce a pre- and postcondition into our example, we see
this behaviour:

\small
\begin{lstlisting}
...
    public inc: int ==> int
    inc(i) == ( total := total + i; return total; )
    pre i < 10
    post total < 20;

traces
    T1:
        let a in set {6, 7, 10} in
            obj.inc(a){1, 5}
\end{lstlisting}

\scriptsize
\lstset{style=tool,language=}
\begin{lstlisting}[escapechar=@]
> runtrace Tester`T1
Generated 15 tests
Test 1 = obj.inc(6)
Result = [6, PASSED]
Test 2 = obj.inc(6); obj.inc(6)
Result = [6, 12, PASSED]
Test 3 = obj.inc(6); obj.inc(6); obj.inc(6)
Result = [6, 12, 18, PASSED]
Test 4 = obj.inc(6); obj.inc(6); obj.inc(6); obj.inc(6)
Result = [6, 12, 18, Error 4072: Postcondition failure: post_inc in 'Counter' at line 15:16, FAILED]
Test 5 = obj.inc(6); obj.inc(6); obj.inc(6); obj.inc(6); obj.inc(6)
Test 5 FILTERED by test 4
Test 6 = obj.inc(7)
Result = [7, PASSED]
Test 7 = obj.inc(7); obj.inc(7)
Result = [7, 14, PASSED]
Test 8 = obj.inc(7); obj.inc(7); obj.inc(7)
Result = [7, 14, Error 4072: Postcondition failure: post_inc in 'Counter' at line 15:16, FAILED]
Test 9 = obj.inc(7); obj.inc(7); obj.inc(7); obj.inc(7)
Test 9 FILTERED by test 8
Test 10 = obj.inc(7); obj.inc(7); obj.inc(7); obj.inc(7); obj.inc(7)
Test 10 FILTERED by test 8
Test 11 = obj.inc(10)
Result = [Error 4071: Precondition failure: pre_inc in 'Counter' at line 14:11, INCONCLUSIVE]
Test 12 = obj.inc(10); obj.inc(10)
Test 12 FILTERED by test 11
Test 13 = obj.inc(10); obj.inc(10); obj.inc(10)
Test 13 FILTERED by test 11
Test 14 = obj.inc(10); obj.inc(10); obj.inc(10); obj.inc(10)
Test 14 FILTERED by test 11
Test 15 = obj.inc(10); obj.inc(10); obj.inc(10); obj.inc(10); obj.inc(10)
Test 15 FILTERED by test 11
Executed in 0.075 secs. 
Some tests failed or indeterminate
\end{lstlisting}
\lstset{style=mystyle}
\lstset{language=VDM++}
\normalsize

The \texttt{inc} operation now has a pre-condition that the argument must be
less than 10 and a postcondition that the resulting total must be less than 20.
The trace makes 1 to 5 calls to the \texttt{inc} operation with arguments 6, 7
and 10, respectively.

The first three tests are fine, but \texttt{Test 4} fails because the fourth call to
\texttt{inc(6)} pushes the total over the limit. This is therefore a
post-condition \texttt{FAILED} test, and the error message is listed along with the
results of the earlier operation calls. Test 5 then tries to do the same, but
adds a further call. This must fail in the same place as \texttt{Test 4}, because\texttt{Test 4}
is the ``stem'' of \texttt{Test 5}. Therefore this test is ``\texttt{FILTERED} by \texttt{Test 4}''.
Similarly, \texttt{Test 8} fails and Tests 9 and 10 are filtered by this failure.

Test 11 fails on the first call to \texttt{inc(10)}, since the argument must be
less than 10. This produces an \texttt{INDETERMINATE} error because we are not
sure whether this is a problem with the trace or the specification being tested.
Lastly, Tests 12 to 15 are filtered by Test 11, since they would behave the same
way.

At the end of the run, the \texttt{runtrace} command indicates that some tests
failed or were indeterminate, just to remind you.

\section{Trace Reduction}

A trace may expand to many millions of tests and these may take many hours to
execute. This can make large traces difficult to work with, both while the trace
is being developed and subsequently when traces are used to check that a change
to a specification is sound.

Therefore the trace system provides the means to reduce the number of tests that
are generated. This limited number can then be checked more quickly, which
naturally does not provide the confidence of a full execution, but is convenient
to work with when a sample of tests is sufficient.

Trace reduction can be achieved in four different ways:

\begin{itemize}
\item \emph{RANDOM} reduction. This is the simplest kind of reduction, and is
used to randomly select tests from the full set. For example, a 1\% RANDOM 
reduction of a trace that expands to a million tests would select 10,000 tests
from the set. The pseudorandom selection is seeded, so that a consistent subset
of the tests can be selected.
\item \emph{SHAPES\_NOVARS} reduction. The ``shape'' of a test means the sequence
of operation names (regardless of arguments passed), and the idea of shaped
reduction is to preserve \emph{at least one test} of every shape in the full
set. So for example, a test that calls \texttt{[opA(1), opB(2), opC(3)]} would
be considered the same shape as a test that calls \texttt{[opA(111), opB(222),
opC(333)]}, but different to a test that calls \texttt{[opX(1), opY(2)]}. So a
shaped reduction of 1\% of one million tests would try to select 10,000 tests,
but it also guarantees to include at least one test of every shape within the
million, even if that means the reduction is (say) 1.5\%.
\item \emph{SHAPES\_VARNAMES} reduction. The simple interpretation of a
shape above only looks at the names of the operations called. This second kind
of shaped reduction looks at the variable names used in the \texttt{let}
bindings and constants as well as the names of the operations. So with this kind
of reduction, \texttt{[opA(a), opA(b)]} is different to \texttt{[opA(x),
opA(y)]}, even if the values of ``a'' and ``x'' can sometimes be the same.
Typically, different variable names are used in different parts of a complex
trace and so this reduction method is trying to select all of the different
parts of a trace, even if the sequence of operation names produced is the same
as another part of the trace.
\item \emph{SHAPES\_VARVALUES} reduction. The third type of shaped reduction is
even more specific about what constitutes a shape, taking into account the
\emph{value} of the variables used as well as their names. So \texttt{[opA(a)]}
will be considered a diffent shape to another \texttt{[opA(a)]} elsewhere, as
long as ``a'' is bound to a different value.
\end{itemize}

The following simple trace, using the \texttt{Counter} example from previous
chapters, illustrates RANDOM reduction:

\small
\begin{lstlisting}
traces
    T1:
        let a in set {1, ..., 1000} in
        let b in set {1, ..., 1000} in
            ( obj.inc(a); obj.dec(b) );
\end{lstlisting}

\lstset{style=tool,language=}
\scriptsize
\begin{lstlisting}[escapechar=@]
> filter
Usage: filter %age | RANDOM | SHAPES_NOVARS | SHAPES_VARNAMES | SHAPES_VARVALUES | NONE
Trace filter currently 100.0% NONE (seed 0)
 
> filter random
Trace filter currently 100.0% RANDOM (seed 0)

> filter 1%
Trace filter currently 1.0% RANDOM (seed 0)

> seedtrace 1234
Trace filter currently 1.0% RANDOM (seed 1234)

> runtrace Tester`T1
Generated 1000000 tests, reduced to 10000, in 1.571 secs. 
Test 66 = obj.inc(1); obj.dec(66)
Result = [1, -65, PASSED]
Test 155 = obj.inc(1); obj.dec(155)
Result = [1, -154, PASSED]
Test 323 = obj.inc(1); obj.dec(323)
Result = [1, -322, PASSED]
Test 419 = obj.inc(1); obj.dec(419)
Result = [1, -418, PASSED]
...
Test 999815 = obj.inc(1000); obj.dec(815)
Result = [1000, 185, PASSED]
Test 999894 = obj.inc(1000); obj.dec(894)
Result = [1000, 106, PASSED]
Test 999992 = obj.inc(1000); obj.dec(992)
Result = [1000, 8, PASSED]
Executed in 4.932 secs. 
\end{lstlisting}
\lstset{style=mystyle}
\lstset{language=VDM++}
\normalsize

The ``filter'' and ``seedtrace'' commands are used to set the filtering
required, then the trace is executed as normal. But we see that the million
tests have been reduced to 10,000 (taking a few seconds). Then instead of trying
every combination of ``a'' and ``b'', the filtering selects them at random,
starting with \texttt{a = 1, b = 65} and ending with \texttt{a = 1000, b = 992}.

The next example illustrates shaped reduction. The trace joins together two
\texttt{let} bindings, repeating each call once and twice. This would normally
produce 36 tests.

\small
\begin{lstlisting}
traces
    T1:
        let a in set {1, 2, 3} in obj.inc(a){1,2};
        let b in set {1, 2, 3} in obj.dec(b){1,2}
\end{lstlisting}

\lstset{style=tool,language=}
\begin{lstlisting}
> filter 1
Trace filter currently 1.0% RANDOM (seed 0)

> filter shapes_novars
Trace filter currently 1.0% SHAPES_NOVARS (seed 0)

> rt Tester`T1
Generated 36 tests, reduced by SHAPES_NOVARS, in 0.001 secs. 
Test 1 = obj.inc(1); obj.dec(1)
Result = [1, 0, PASSED]
Test 2 = obj.inc(1); obj.inc(1); obj.dec(1)
Result = [1, 2, 1, PASSED]
Test 7 = obj.inc(1); obj.dec(1); obj.dec(1)
Result = [1, 0, -1, PASSED]
Test 8 = obj.inc(1); obj.inc(1); obj.dec(1); obj.dec(1)
Result = [1, 2, 1, 0, PASSED]
Executed in 0.006 secs. 
All tests passed
>
\end{lstlisting}
\lstset{style=mystyle}
\lstset{language=VDM++}
\normalsize

Here, we try to reduce this to 1\%, but using the SHAPES\_NOVARS option. A
reduction of 1\% of 36 tests ought to produce a single test (reduction will
never produce zero tests), but in fact the shaped reduction produces four. You
can see that there is one case of each ``shape'': one inc and one dec; two incs
and one dec, and so on. So the idea is that the shaped reduction has given a
representative sample of all of the possible shapes, disregarding variable names
or values.

Reducing this trace using SHAPES\_VARNAMES produces the same number of shapes,
since the two calls always use the same variable names. But reduction using
SHAPES\_VARVALUES regards all of the tests as different shapes, because the
variable/value/operation combinations are different in all of the tests.

In practice, the most useful reductions are RANDOM and SHAPES\_NOVARS. Random
reductions give a simple way to cut down a large number of tests. Shaped
reduction is choosing one example of every ``path'' that the trace is taking the
specification through, which is often closely related to the different use cases
that the system has.


\section{How Does Trace Expansion Work?}

The sections above have given an overview of all the trace operators, and there
are some examples of combinations of operators. But to see how traces are
expanded in general, we need to look at traces from a different point of view.
The syntax of traces is deliberately made similar to the syntax of VDM-SL, but
to understand how operators combine to produce multiple tests, it helps to look
at operators as though they followed a separate ``expansion'' grammar. In the
description that follows, a \texttt{set} is a set of tests:

\begin{itemize}
\item \texttt{set = object.opname(args)}. The simplest form of a trace is a set
that comprises a single call to an operation or function with arguments. The
arguments can be symbolic, and bound to various values by the \texttt{let}
operator described below.
\item \texttt{set = set1; set2; ...; setn}. A set of tests may be formed from
an ordered sequence of sets. This expands to all possible selections of one test
from each of the sets. In its simplest form, this could be a sequence of
operation calls which therefore just expands to one test. But a combination of
sets of tests results in a set of the product of the sizes of those sets.
\item \texttt{set = set1 ?}. A set of tests may be formed from another set with
a \texttt{?} operator. This produces the same set, but includes a ``skip'' step.
\item \texttt{set = set1 \{n[, n]\}}. A set of tests may be formed from another
set with a \texttt{\{n\}} or \texttt{\{n1, n2\}} operator. This produces a set
with every member of the original set repeated \texttt{n} times, or between
\texttt{n1} and \texttt{n2} times (inclusive).
\item \texttt{set = set1 *|+}. A set of tests may be formed from another set
with a \texttt{*} or \texttt{+} operator. This produces another set with every member
of the original set repeated from 0 to N times (with \texttt{*}) or 1 to N times
(with \texttt{+}). The value of N is tool dependent, but defaults to 5.
\item \texttt{set = set1 | set2 | ... | setn}. A set of tests may be formed by
combining a number of other test sets with a \texttt{|} operator. This produces
a set with the union of the other sets.
\item \texttt{set = || (set1, set2, ..., setn)}. A set may be formed from the
permutations of a number of other sets. This produces a set with each
permutation of each selection of one test from each set.
\item \texttt{set = let <multiple bind> [be st <cond>] in set1}. A set of
tests may be formed from a \texttt{multiple bind}, which expands to the
substitution of all the possible the bound values in the original set.
\item \texttt{set = let <name> = <exp> [, <name2> = <exp2>, ...] in set1}. A set
of tests may be evaluated in a scope that defines name/value pairs. This does
not increase the number of tests in the set, but just binds free variables.
\end{itemize}

For example, if (for brevity) we say that a test with a single call to
\texttt{obj.opA()} is written as ``[A]'', and similarly ``[B]'' and ``[C]'' 
for other operation calls, and ``[-]'' for a skip, then we can say the 
following trace operators produce these sets of tests:

\small
\lstset{style=tool,language=}
\begin{lstlisting}[escapechar=@]
A? = { [A], [-] }
A;B = { [AB] }
A;B? = { [AB], [A] }
A* = { [-], [A], [AA], [AAA], [AAAA], [AAAAA], ... }
A+ = { [A], [AA], [AAA], [AAAA], [AAAAA], ... }
A{3} = { [AAA] }
A{1,3} = { [A], [AA], [AAA] }
A | B = { [A], [B] }
A | B? = { [A], [B], [-] )
|| (A, B, C) = { [ABC], [ACB], [BAC], [BCA], [CAB], [CBA] }
|| (A, (B;C)) = { [ABC], [BCA] }
|| (A, B+) = { [AB], [BA], [ABB], [BBA], [ABBB], [BBBA], ... }
let a in set {1,2,3} in A(a) = { [A(1)], [A(2)], [A(3)] }
let b : bool * bool in B(b) = {
    [B(mk_(true, true))], [B(mk_(true, false))],
    [B(mk_(false, true))], [B(mk_(false, false))]
}
let z = 1 in B(z) = { [B(1)] }
\end{lstlisting}
\lstset{style=mystyle}
\lstset{language=VDM++}
\normalsize

Note that the repeat limits in a trace (like \{1,3\}) must be numeric literals.
But values in a multiple bind set or sequence can be variables, either bound
earlier in the trace or other fields within scope of the trace inside the object
or module where it is defined. Similarly, the values in the right hand side of
let definitions can be variables within the trace or the object/module scope.

\section{Language Considerations}

Combinatorial tests are available for both VDM-SL and VDM++/VDM-RT. The process 
of trace expansion and execution is very similar in all cases, but there are
some differences that are described below.

\subsection{Traces in VDM-SL}

Traces are added in a ``traces'' section within a VDM-SL specification. This can
either be within one or more modules or within a flat specification. The name of
the traces in a module are implicitly exported, so they are referred to as
\texttt{<modulename>`<tracename>}. You can omit the module name if it is the
default module.

The VDM-SL specification that is equivalent to the example used above is like
this:

\small
\begin{lstlisting}
module Counter
exports all
definitions

state S of
    total:int
init s == s = mk_S(0)
end

operations
    inc: int ==> int
    inc(i) == ( total := total + i; return total; );

    dec: int ==> int
    dec(i) == ( total := total - i; return total; )
end Counter

module Tester
imports from Counter all
definitions

traces
    T1: Counter`inc(1)*;

end Tester
\end{lstlisting}
\normalsize

And in the VDM-SL command line, that would be executed as follows. Note that
Tester is not the default module, so the trace name is qualified:

\scriptsize
\lstset{style=tool,language=}
\begin{lstlisting}[escapechar=@]
> modules
Counter (default)
Tester
> runtrace Tester`T1
Generated 6 tests in 0.002 secs. 
Test 1 = skip
Result = [(), PASSED]
Test 2 = inc(1)
Result = [1, PASSED]
Test 3 = inc(1); inc(1)
Result = [1, 2, PASSED]
Test 4 = inc(1); inc(1); inc(1)
Result = [1, 2, 3, PASSED]
Test 5 = inc(1); inc(1); inc(1); inc(1)
Result = [1, 2, 3, 4, PASSED]
Test 6 = inc(1); inc(1); inc(1); inc(1); inc(1)
Result = [1, 2, 3, 4, 5, PASSED]
Executed in 0.019 secs. 
All tests passed
>
\end{lstlisting}
\lstset{style=mystyle}
\lstset{language=VDM++}
\normalsize

This trace is very similar to the VDM++ example. The \texttt{Counter} module has a
single state that is equivalent to the VDM++ ``total'' instance variable. Note
that this is reset to zero automatically before each test is executed. This is
because each test re-initializes the specification, and the module
state has an ``\texttt{\textbf{init}}'' clause that sets the total to zero.

Notice also that the operation calls are not applied to a Counter object, unlike
VDM++.

\subsection{Traces in VDM++ and VDM-RT}

Traces are added in a ``traces'' section within a VDM++ or VDM-RT
specification, inside one or more classes. In effect, the name of the trace is
a public static symbol, so it is referred to as
\texttt{<classname>`<tracename>}, as we have seen in the examples above. You
can omit the class name if that is the default class.

Although a VDM++ trace is effectively a static scope, and can call static
operations directly (similar to a VDM-SL trace), every test execution occurs in
a \emph{new instance} of the containing class - in our examples, in a new Tester
instance. This means that objects created within the Tester's construction will
be freshly initialized and ready for use in each test run. In the example, the
Counter object \texttt{obj} is created for each test, because the instance
variable is initialized at construction.

\subsection{Expansion and Execution Considerations}

The process of running a combinatorial test has two phases: expanding the trace
to a number of test definitions; and subsequently executing those definitions.
The trace expansion typically does not take very long, since it is only
constructing a tree of iterators that are capable of generating the tests one
after another. The subsequent execution of those tests can obviously take a long
time, depending on how many there are.

We have seen (above) how the specification is initialized before test execution,
and the state of the module or class is available to the trace, but care must be
taken if operations or functions within the environment are used as part of a
trace. This is because some expressions are evaluated during trace expansion and
some during test execution. For example:

\small
\begin{lstlisting}
...
functions
    private static range: int * int -> set of int
    range(a, b) == {a, ..., b};

values
    Z = 100;

traces
    T1: let x in set range(3, 5) in
        let y in set range(x, x+2) in
            obj.inc(Z + x + y);
\end{lstlisting}
\normalsize

In this case, the \texttt{range} function is used to create a set for the
multi-binds, and this is executed during \emph{expansion}, once for ``x'' and three
times for ``y''. Similarly, the right hand side of simple let definitions are
executed during trace expansion. But the addition of \texttt{Z + x + y} in the
argument to \texttt{inc} is called during \emph{execution} (once for each test).

Trace generation starts inside a fresh object instance of the class (or
initialized module) that contains the trace. So if operations are called during
trace \emph{expansion}, these can modify state and so affect subsequent
operation calls elsewhere in the expansion. This can become very confusing, and
it is not a recommended trace design strategy! On the other hand, calling
functions as part of the trace expansion can make traces easier to understand
and can provide the means to build complex sets that would be difficult to
construct directly within the trace statements.

When a test is listed in the trace output, the arguments that are passed to
operation calls are shown as literals, if possible. As seen in the examples
here, a call to \texttt{obj.inc(1)} is shown, rather than \texttt{obj.inc(a)}. This is
possible whenever arguments can be easily evaluated. For example the \texttt{Z
+ x + y} case above would produce \texttt{Test 1 = obj.inc(106)}, which is 100 +
3 + 3. But if the argument is a more complex expression involving operation
applies or new object creation, these cannot be evaluated and so the argument
expression is listed ``as is''. For example, if the trace above is changed to
call \texttt{obj.inc(max(x, y))}, the test would be listed with ``\texttt{x}''
and ``\texttt{y}'' rather than their current values:

\small
\lstset{style=tool,language=}
\begin{lstlisting}[escapechar=@]
> runtrace Tester`T1
Generated 9 tests in 0.006 secs. 
Test 1 = inc(max(x, y))
Result = [3, PASSED]
Test 2 = inc(max(x, y))
Result = [3, PASSED]
...
\end{lstlisting}
\lstset{style=mystyle}
\lstset{language=VDM++}
\normalsize

\chapter{Combinatorial Testing Patterns}
\label{chap:patterns}

\section{Traces and Test Operations}

After expanding a trace to all of the call sequences that match, a
test is ultimately just a sequence of operation or function calls. But these
calls do not have to be the primary operations that drive the specification. In
some cases, it makes sense for traces to expand to a set of tests that call
\emph{test operations} that maintain their own state and exercise the main
specification, checking the responses and the main state. Checks like this, that
do not directly form part of the constraints of the specification, are usually
called \emph{validation conjectures}. The task of the testing operations is
therefore to check that, whatever the call sequence made by the trace, the
validation conjectures for the specification are maintained.

For example, a specification may describe how a sequence of parts are produced
as calls are made to \texttt{newPart(), adjustPart(), completePart()}. The
process of creation of individual parts may well involve preconditions,
postconditions, type invariants and so on. But there may also be a requirement
that (say) over time, the total number of parts of type A and type B never
differ by more than a tolerance. This is a validation conjecture: it is not
directly stated in the constraints of the specification, but it is a behaviour
that must be manifest by the system over time. Therefore the trace(s) for such
a system can call the main operations via test operations, like
\texttt{testNewPart()} and so on, and the \texttt{testCompletePart()} operation
can check the history of parts created to validate that the tolerance is always
respected, regardless of the sequence of operations that the trace tries. Note
that these test operations can maintain their own state that is private and
separate from that of the main specification.

In such cases, it makes sense to add all of the test operations to a separate
class or module, to make it clear that they are not part of the main
specification. This is illustrated in the Basket Service example in Chapter
\ref{chap:examples}.

\section{Common Patterns}

Experience has shown that the testing of many specifications with traces
requires several common ``patterns'' to create data selections or sequences of
operations. These are presented in this chapter.

\subsection{Sets for ``let'' bindings}

The set of values that is used by a \texttt{let multiple-bind} is usually
shown as an enumeration of literals in examples. But the set value can use any
VDM expression that yields a set. The following cases are generally useful:

\begin{itemize}
  \item Set comprehensions can be used to select values from a larger set that
  meet the membership predicate. This is similar to the use of the \texttt{be
  st} clause, but the filter acts on the members of the set rather than the
  values that are bound:
\small
\begin{lstlisting}
let pair in set {mk_(a, b) | a, b in set VALUES & a > b} in ...
\end{lstlisting}
\normalsize
  \item The \texttt{power} operator can be used to produce all possible subsets
  from another set, though you often have to eliminate the empty set, which is
  produced by the operator:
\small
\begin{lstlisting}
let options in set power ALL_OPTIONS \ {{}} in ...
\end{lstlisting}
\normalsize
  \item If you are using finite types (i.e. types that have a finite number of
  values), then you can use a multiple type bind to conveniently choose all of
  the possible values of that type:
\small
\begin{lstlisting}
let p1, p2 : Product in ...
\end{lstlisting}
\normalsize
  \item Often elements have to be chosen from a set such that they are different
  to previous selections from the same set. This can sometimes be done in a
  single \texttt{let} bind by using a \texttt{be st} clause that states that the
  values are different. But a common usage is to make a selection from a set
  that has had the first choice(s) eliminated:
\small
\begin{lstlisting}
let a, b in set S be st a <> b in ...
let c in set S \ {a, b} in ...
let d in set S be st d not in set {a, b, c} in ...
\end{lstlisting}
\normalsize
  \item A common requirement is to select permutations of a set. This can be
  done by using the looseness of a set bind. Note the cardinality check in the
  k-permutation example (k=3), and the use of the set pattern and the
  set-of-set \texttt{\{S\}} in the full permutation example - in this case you
  have to know the size of the set to match the pattern:
\small
\begin{lstlisting}
-- k-permutations of 3 values from S
let c1, c2, c3 in set S be st card {c1, c2, c3} = 3 in ...
-- Permutations of all values from S
let {p1, p2, p3, p4, p5} in set {S} in ...
\end{lstlisting}
\normalsize
  \item Something about combinations!!
\end{itemize}

\subsection{Bracketing operations with headers and trailers}

A very common requirement is to make a number of fixed setup calls,
followed by a large number of different test calls, followed by a number of
fixed closedown calls. This pattern emerges naturally from a semi-colon
separated trace sequence with a complex ``middle'':
\scriptsize
\begin{lstlisting}
traces
T:
    SetSystemDate(221220);
    LoadCertificates(RefData);
        test1() |
        test2() |
        test3() |
        test4() |
        test5() |
        test6() |
        test7();
    EndTransaction();
\end{lstlisting}
\normalsize
This produces a set of seven tests, calling tests 1 to 7, each of which is
sandwiched by calls that set up the system and close down the transation.
Clearly the body containing the alternative tests can be arbitrarily
complicated, for example using variable binds to generate many possiblities.

\subsection{Graph searching by traces}

Traces that explore all of the possibile uses cases of a system frequently
need to search a graph that describes the possible paths through the use case.
This is translated into a trace by generating a call for each step through the
graph, where the trace expands to the possibilities at each step. For example:
\scriptsize
\begin{lstlisting}
values
    TESTSTATES : map TestState to map Event to TestState =
    {
        <READY>    |->
        {
            <SEND>        |-> <SENT>
        },

        <SENT> |->
        {
            <RX_NACK>    |-> <RESEND1>,
            <RX_ACK>    |-> <END>,
            <TIMEOUT>    |-> <RESEND1>,
            <PARTIAL>    |-> <SENT>
        },
        ...
    }

traces
    AllTransitions:
        let s1 = <READY> in
        let ev1 in set dom TESTSTATES(s1) in
        let s2 = TESTSTATES(s1)(ev1) in
        let ev2 in set dom TESTSTATES(s2) in
        let s3 = TESTSTATES(s2)(ev2) in
        let ev3 in set dom TESTSTATES(s3) in
        let s4 = TESTSTATES(s3)(ev3) in
        let ev4 in set dom TESTSTATES(s4) in
        let s5 = TESTSTATES(s4)(ev4) in
        let ev5 in set dom TESTSTATES(s5) in
            execute([ev1, ev2, ev3, ev4, ev5]);
 
\end{lstlisting}
\normalsize
Here the system starts in state \texttt{s1}, which must be \texttt{READY}. From
there, a number of events are possible given by the domain of a lookup of the
state in the \texttt{TESTSTATES} map. One of these events is selected as
\texttt{ev1}, which then moves us to state \texttt{s2}, and so on. After five
events have been generated, they are passed to an \texttt{execute} method which
uses the events to test that particular path through the graph.

Note that a graph searching cascade like this can generate tens of thousands of
possibilities very quickly, even with comparatively simple graphs.

\subsection{Building data in stages}

A cascade of let bindings can be used to build a complex data structure in
stages, rather than each binding having to create a complete value. For example:
\scriptsize
\begin{lstlisting}
EPATest:
    let a, b, c, d, e in set    -- Pick five branches
    {
        mk_B([1],        0),
        mk_B([-1, 2],    1),
        mk_B([1, 2, 0],  2)
    }
    in

    let branches in set         -- Between one and five of them
    {
        [a],
        [a,b],
        [a,b,c],
        [a,b,c,d],
        [a,b,c,d,e]
    }
    in
 
    let epa = mk_InputFile(     -- EPA InputFile from those B values
        mk_Header(),
        mk_PaymentSummary(getPS(branches)),
        {
            mk_Branch(
                mk_SummaryOfCharge(mk_token(mk_("MID", B)), ...),
                {
                    mk_RecordOfCharge(... mk_("TxnNum", B, i) ...)
                    | i in set inds branches(B).ROCs
                },
                {
                    mk_Adjustment(mk_token(mk_(B, i)))
                    | i in set {1, ..., branches(B).ADJs}
                }
            )
            | B in set inds branches
        },
        mk_Trailer(getTR(branches) + 1)
    )
    in

    -- Finally, transform the EPA file into the various output formats.
    (
        transformAudit(epa);
        transformC4D(epa)
    );
\end{lstlisting}
\normalsize
This example creates \texttt{InputFiles}, which contain various records that
relate to electronic payments taken from a branch of a business. The objective of the
trace is to check a large number of different input files with different numbers
of branches and various record types.

Creating such files in a single trace step would be difficult, if not
impossible. But here, we see that the generation starts with a selection of
``B'' values from a set of possibilities. Then sequences of between one and five
of these values is created. Then these ``branch'' sequences are used in nested
set comprehensions to create an \texttt{InputFile}. Lastly, every InputFile
created is processed by a couple of operations.

\subsection{Oracle functions}

In many cases, the expected result of an operation or function call is too
complicated to predict simply in a trace - assuming you have a support function
that checks the result with (say) a postcondition:

\scriptsize
\begin{lstlisting}
    assert: seq1 of nat * seq1 of nat +> bool
    assert(data, expected) == data = expected
    post RESULT = true;
\end{lstlisting}
\normalsize

Here, we assume that the ``data'' argument passed comes from some processing in
the specification, but where does the ``expected'' value come from? It could be
a literal in the trace, but this is not practical for non-trivial examples, so
it is common to create support functions that produce answers that are correct
\emph{by definition}. Such functions are called \emph{oracles}. For example:

\scriptsize
\begin{lstlisting}
Scenarios:
    let s = mk_Service(...) in
    let q in set ... in
    let mk_(first, second) in set ... in 
    let test = mk_Test(
        [s],
        keyStrokesFor(q, first, second),    -- Expected keystrokes
        basketFor(first, second))           -- Expected basket
    in
        run(test);
\end{lstlisting}
\normalsize

In this example, a \texttt{Test} record is created that includes a Service, the
expected keystrokes and the expected basket result for a retail application in
a given scenario (use case). The \texttt{run} operation uses the key strokes to
drive the specification under test and then check that the expected basket is
produced correctly. Both of these functions are oracles.

\subsection{Using positive and negative sense checks}

The natural way to think about testing a specification is to consider all of the
success paths and then design traces that exercise those paths. But in many
cases, a specification is also required to ``fail'' in the correct way; that is,
there certain inputs or call sequences that require a specific error condition
to be generated, even though that is a failure in some sense.

There is an example of this in the Luhn specification described in Chapter
\ref{chap:examples}. The Luhn algorithm is a kind of checksum. Therefore it is
\emph{required} to fail if a piece of data is corrupted in specific ways, which
shows that the checksum is doing its job, detecting the corruption. So the Luhn
specification tests deliberately corrupt a piece of data and then verify that
the algorithm generates a \emph{different} check digit - i.e. they verify that a
check with the original check digit fails correctly:

\scriptsize
\begin{lstlisting}
    checkFail: seq1 of nat * nat * nat ==> bool
    checkFail(data, expected, base) ==
        return luhn(data, base) <> expected    -- Expect failure
    post RESULT = true;
\end{lstlisting}
\normalsize

\section{Avoiding Test Explosions}

It is extremely easy to write a comparatively simple looking trace definition
that expands into a collection of tests that is so large that it is not
practical to execute, either because it would take too long or because the
generation process takes up too much memory.

The following tips will help you to avoid this pitfall:

\begin{itemize}
  \item \emph{Start small}. It is tempting to write traces as clearly as
  possible to start with, and that may lead to the binding of values from sets
  of data or types with many values. The combinatorial expansion process will
  then either multiply these data sizes together, or in some cases generate 
  combinations that depend on the product of the \emph{factorial} of the data
  sizes. So a modest set of 50 values might therefore generate of the order of
  $10^{64}$ tests (i.e. the factorial of 50). So start small: design and test 
  traces with small example sets and types, and only expand the data selections
  when you can see that the trace is expanding as you require.
  \item \emph{Split up traces}. The alternation operator, \texttt{|}, will join
  together two sets of tests in a trace. It may therefore be tempting to write a
  single trace that is composed of many parts, testing different parts of the
  system, joined by alternation. This has the advantage that the whole
  specification can be tested by running one trace. But it also means that the
  trace expands to the sum of all of the tests within all of the parts. That is
  much better than a product or factorial combination, but if the parts are
  genuinely separate, you will be able to do more tests in the parts by
  separating them into multiple traces.
  \item \emph{Be careful with multiple-binds and power sets}. As mentioned
  above, factorial scaling is extremely expensive. This occurs most commonly
  with multiple-binds that are used for k-permutations, and with the
  \texttt{power} operator on sets.
\end{itemize}

\chapter{Combinatorial Testing Examples}
\label{chap:examples}

\section{The LUHN Check Digit Model}
\section{The Basket Service Model}

\appendix
\chapter{Combinatorial Testing Syntax}
\label{chap:CTsyntax}

\Ruledef{traces definitions}{\Lop{traces},
   \OptPt{\Ruleref{named trace}, \SeqPt{\Lit{;}, \Ruleref{named trace}}}
}

\Ruledef{named trace}{\Ruleref{identifier},
  \SeqPt{\Lit{/}, \Ruleref{identifier}},
  \Lit{:},\Ruleref{trace definition list}
}

\Ruledef{trace definition list}{
  \Ruleref{trace definition term},
  \SeqPt{\Lit{;}, \Ruleref{trace definition term}}
}

\Ruledef{trace definition term}{
  \Ruleref{trace definition}, \SeqPt{\Lit{|}, \Ruleref{trace definition}}
}

\Rule{trace definition}{
  \Ruleref{trace binding definition} \dsep
  \Ruleref{trace repeat definition}
}

\Ruledef{trace binding definition}{
  \Ruleref{trace let def binding} \dsep
  \Ruleref{trace let best binding}
}

\Ruledef{trace let def binding}{
  \Lop{let}, \Ruleref{local definition},
  \SeqPt{ \Lit{,}, \Ruleref{local definition}}, \lfeed
  \Lop{in}, \Ruleref{trace definition}
}

\Ruledef{trace let best binding}{
  \Lop{let}, \Ruleref{multiple bind},
  \OptPt{\Lop{be}, \Lop{st}, \Ruleref{expression}}, \lfeed
  \Lop{in}, \Ruleref{trace definition}
}

\Ruledef{trace repeat definition}{
  \Ruleref{trace core definition},
  \OptPt{\Ruleref{trace repeat pattern}}
}

\Ruledef{trace repeat pattern}{
  \Lit{*} \dsep
  \Lit{+} \dsep
  \Lit{?} \dsep
  \Lit{\{}, \Ruleref{numeric literal}, \OptPt{\Lit{,}, \Ruleref{numeric literal}, \Lit{\}}}
}

\Rule{trace core definition}{
  \Ruleref{trace apply expression} \dsep
  \Ruleref{trace concurrent expression} \dsep
  \Ruleref{trace bracketed expression}
}

\Rule{trace apply expression}{
  \Ruleref{call statement}
}

\Rule{trace concurrent expression}{
  \Lit{||}, \Lit{(}, \Ruleref{trace definition}, \lfeed
  \Lit{,}, \Ruleref{trace definition}, \lfeed
  \SeqPt{\Lit{,}, \Ruleref{trace definition}}, \Lit{)}
}

\Ruledef{trace bracketed expression}{
  \Lit{(}, \Ruleref{trace definition list}, \Lit{)}
}

\chapter{Overture Screenshots}
\label{chap:OvertureShots}

\chapter{Example Model Listings}
\label{chap:listings}
The VDM models listed in this appendix are described in detail in Chapter
\ref{chap:examples}

\section{The LUHN Check Digit Model}
\section{The Basket Service Model}

\bibliographystyle{newalpha}

\bibliography{../bib/dan}

\end{document}


